\section{Welfare}\label{sec:welfare}

The price effect is one component of welfare cost; two others are routinely omitted in reduced-form work. Allocative deadweight loss arises because the set-aside routes the contract to an SME winner whose cost may exceed the lowest cost in the unrestricted pool. Tax-distortion loss arises because the additional government outlay is financed from distortionary taxation; each marginal real raised costs $(1+\lambda)$ in shadow-price terms. Transfer between the government and SME bidders is, by contrast, zero-sum in welfare---it shifts surplus across agents but does not destroy it. The decomposition I report isolates the two destruction components and leaves the transfer implicit.

Formally, per-auction welfare loss under the set-aside regime is
\begin{equation}
  \text{Loss} = \underbrace{c_{(1)}^{S_3} - c_{(1)}^{S_1}}_{\mathrm{DWL}_{\text{alloc}}}
              \,+\, \underbrace{\lambda \cdot (p^{S_3} - p^{S_1})}_{\mathrm{MCPF\ dist.}}.
\end{equation}
The first term is the allocative wedge: the extra production cost of serving the auction under the restricted regime. The second is the tax distortion on the government's extra outlay $\Delta_{\text{gov}} = p^{S_3} - p^{S_1}$. Both are reported in ref-price units and as a percentage of the open-regime price $p^{S_1}$.

\subsection{Point estimates}

Table~\ref{tab:v3_welfare} shows the decomposition at $\lambda = 0.30$, the \citet{ballard1985} benchmark value adopted for Brazil in subsequent literature. In non-pharma, $\Delta_{\text{gov}} = 0.24$, $\mathrm{DWL}_{\text{alloc}} = 0.15$ (62\% of $\Delta_{\text{gov}}$), MCPF distortion $= 0.073$, and total welfare loss is $0.223$---that is, \textbf{28.7\% of the open-regime procurement price $p^{S_1}$}. In pharma the numbers are larger: $\Delta_{\text{gov}} = 0.31$, $\mathrm{DWL}_{\text{alloc}} = 0.21$, MCPF distortion $= 0.092$, total loss $0.298 = $ \textbf{45.5\% of $p^{S_1}$}. The implicit transfer to SMEs is $\Delta_{\text{gov}} - \mathrm{DWL}_{\text{alloc}} = 0.09$ (non-pharma) and $0.10$ (pharma) in ref-price units.

\subsection{Confidence intervals}

Cluster bootstrap at the auction level (B=500, Table~\ref{tab:v3_welfare_ci}) delivers 95\% CIs of $[21.3\%, 35.3\%]$ in non-pharma and $[35.1\%, 56.2\%]$ in pharma at $\lambda = 0.30$. Neither interval crosses the v2 reduced-form estimate of 11.8\%; v3 is comfortably larger in both classes. The CI widths of 14--21 percentage points are substantive relative to the point estimates but leave no ambiguity about the direction of the departure from v2.

\subsection{Sensitivity to $\lambda$}

For $\lambda \in \{0.20, 0.30, 0.40\}$---the range that brackets the benchmark 0.30 value of \citet{ballard1985} and its subsequent refinements for emerging economies---the welfare loss in non-pharma runs 25.6\%, 28.7\%, 31.8\%; in pharma, 40.8\%, 45.5\%, 50.2\%. Each 0.10 increment in $\lambda$ adds roughly 3--5 percentage points to the total. The qualitative result is robust to the MCPF choice within the reasonable range.

\subsection{Why v3 is two-to-four times v2}

Two mechanical changes drive the v3 increase:
\begin{enumerate}
\item \textbf{UH-clean $F_c$ has a heavier right tail.} The losers-only ECDF used in v2 omits the lowest cost per auction. The BLP shrinkage applied to all bidders in v3 replaces that with a tail that includes near-winners. This tail drives $\mathrm{DWL}_{\text{alloc}}$ upward by 30--50\% relative to v2.
\item \textbf{MCPF is made explicit.} v2 did not price the tax distortion on the government's extra payment; it reported only the allocative wedge. At $\lambda = 0.30$, MCPF adds 7--9 percentage points to the total loss.
\end{enumerate}
The 28--46\% range is closer to \citet{krasnokutskaya2011}'s estimates for California highways (20--30\% welfare cost of the 10\% price preference they study) than to the 7--11\% benchmarks from US state-level procurement \citep{marion2007,athey2013}. The difference plausibly reflects the more binding nature of a full set-aside relative to a price preference: a set-aside truncates the pool entirely, which generates larger allocative losses even before MCPF is priced.
